\documentclass{article}
\usepackage{hyperref}
\usepackage{Style}

\nocite{*} % Comentar si quiero citar
%\addbibresource{bibliografia.bib} % Quitar el comentado si quiero usar bibliografia

\begin{document}

\begin{minipage}{2.5cm}
    \includegraphics[width=2cm]{imagen_puc.jpg}
\end{minipage}
\begin{minipage}{12cm}
    {\sc Pontificia Universidad Católica de Chile\\
    Facultad de Matemáticas\\
    Departamento de Matemática\\
    Profesor: Manuel Cabezas -- Ayudante: Benjamín Mateluna}
\end{minipage}
\vspace{2ex}

{\centerline{\bf Introducción a la Geometría - MAT1304}
\centerline{\bf Ayudantía 12}}
\centerline{\bf 21 de abril de 2026}

\vspace{1cm}
\noindent\textbf{Problema 1.} Sea $x\in(0,\frac{\pi}{2})$ tal que $sec(x)-tan(x)=2$. Calcule
$sec(x)+tan(x)$.

\vspace{5mm}
\noindent\textbf{Problema 2.} Verifique las identidades:
\begin{equation*}
    (a)\hspace{1mm}\frac{sen(x)}{1-cos(x)}=csc(x)+cot(x)\hspace{1cm}
    (b)\hspace{1mm}tan(x)+cot(x)=sec(x)\cdot csc(x)\hspace{1cm}
    (c)\hspace{1mm}cos\left(\frac{\pi}{2}-x\right)=sen(x)
\end{equation*}

\vspace{5mm}
\noindent\textbf{Problema 3.} Para $x\in\R$, demuestre que
\begin{equation*}
    cos^{2}(x)=\frac{cos(2x)+1}{2}\hhtext{y}
    sen^{2}(x)=\frac{1-cos(2x)}{2}
\end{equation*}

\vspace{5mm}
\noindent\textbf{Problema 4.} Demuestre que para todo $n\in\N$ y para todo $\alpha\in\R$ tal que
$sen(\alpha)\neq0$ se tiene que
\begin{equation*}
    cos(\alpha)\cdot cos(2\alpha)\cdots cos(2^{n}\alpha)
    =\frac{sen(2^{n+1}\alpha)}{2^{n+1}sen(\alpha)}
\end{equation*}

\vspace{1cm}
\noindent\textbf{\underline{Ejercicios Propuestos:}}

\vspace{5mm}
\noindent\textbf{Extra 1.} Exprese el área de $\triangle ABC$, sabiendo el largo del lado 
$\overline{BC}$, $\overline{AB}$ y el valor del ángulo $\angle ABC$.

\vspace{5mm}
\noindent\textbf{Extra 2.} Encuentre un valor explícito para $sen(\frac{\pi}{12})$ y
$cos(\frac{\pi}{12})$. Calcule
\begin{equation*}
    cos^{4}\left(\frac{\pi}{24}\right)-sen^{4}\left(\frac{\pi}{24}\right)
\end{equation*}

%\printbibliography % Quitar el comentado si quiero usar bibliografia

\end{document}
